%% ECON 282E -- Session 5, Deck B1: Projection -- Making a Residual Small
%% Source: Quant_Macro Lec.9 frames 37-43 (L969-1137).  That block follows
%% Fernandez-Villaverde & Guerron-Quintana, Solution Methods IV, closely; it is
%% rewritten here and cited.  Source defects are recorded in slides/SOURCES.md,
%% not on the slides.
%% Numbers and coordinates from tools/figures/s05_numbers.json, key
%% "projection_residual", "fig_residual".

%% ECON 282E -- Foundations of Macroeconomics (UC Riverside, Fall 2026)
%% Shared Beamer preamble for every deck in slides/S01 ... slides/S10.
%%
%% Usage (a deck lives in slides/S0X/ and is compiled from that folder):
%%   %% ECON 282E -- Foundations of Macroeconomics (UC Riverside, Fall 2026)
%% Shared Beamer preamble for every deck in slides/S01 ... slides/S10.
%%
%% Usage (a deck lives in slides/S0X/ and is compiled from that folder):
%%   %% ECON 282E -- Foundations of Macroeconomics (UC Riverside, Fall 2026)
%% Shared Beamer preamble for every deck in slides/S01 ... slides/S10.
%%
%% Usage (a deck lives in slides/S0X/ and is compiled from that folder):
%%   \input{../preamble/course_preamble.tex}
%%   \decktitle{1}{A1}{Who I Am \& Why This Course}{September 24, 2026}
%%   \begin{document}
%%   \makedecktitle
%%   ...
%%
%% Adapted from Courses/AI-ECON-2026/source/shared_preamble.tex (Metropolis theme,
%% concept/caution/example/takeaway boxes, \sectiondivider). Changes for this course:
%%   * course title block (\decktitle / \makedecktitle) instead of \makebeamertitle
%%   * \zh{...} typesets Chinese (book titles) under pdflatex via CJKutf8
%%   * researchbox for the "Research in the wild" frames
%%   * section pages off; TikZ libraries and the course map loaded here
%% Build with pdflatex only (two passes); tools/build_slides.py does this for you.

\documentclass[10pt,english,aspectratio=169]{beamer}

\usepackage[T1]{fontenc}
\usepackage{textcomp}
\usepackage[utf8]{inputenc}
\usepackage{babel}
\usepackage{CJKutf8}

\usepackage{amsmath, amssymb, amsthm, graphicx}
\usepackage{booktabs, multirow, tabularx}
\usepackage{tikz}
\usetikzlibrary{positioning,arrows.meta,shapes,calc,fit,backgrounds}
\usepackage{pgfplots}
\pgfplotsset{compat=1.16}
\usepackage[most]{tcolorbox}
\tcbuselibrary{skins, breakable}
\usepackage{listings}
\usepackage{xcolor}

\lstset{
  basicstyle=\ttfamily\small,
  keywordstyle=\color{blue!80!black}\bfseries,
  commentstyle=\color{green!50!black}\itshape,
  stringstyle=\color{orange!80!black},
  backgroundcolor=\color{gray!8},
  frame=single,
  rulecolor=\color{gray!40},
  breaklines=true,
  showstringspaces=false,
  numbers=none,
}

\usetheme{metropolis}
\usefonttheme{professionalfonts}
\metroset{sectionpage=none}

\definecolor{LightGray}{RGB}{230,230,230}
\definecolor{RedTitle}{RGB}{0,0,200}
\definecolor{VeryLightGray}{RGB}{245,245,245}
\definecolor{DarkText}{RGB}{0,0,0}
\definecolor{UCRBlue}{RGB}{0,45,106}
\definecolor{UCRGold}{RGB}{241,171,0}

% Stage colours of the course arc (used by the course map and elsewhere)
\colorlet{StageTrunk}{blue!12}
\colorlet{StageBranch}{cyan!14}
\colorlet{StageMachine}{gray!18}
\colorlet{StageWall}{red!14}
\colorlet{StageTools}{orange!20}
\colorlet{StageData}{green!16}
\colorlet{StageAgents}{violet!14}

\hypersetup{bookmarks=false,breaklinks=true,pdfborder={0 0 0},colorlinks=true,
  linkcolor=DarkText,urlcolor=RedTitle}

\setbeamercolor{background canvas}{bg=white}
\setbeamercolor{title}{fg=RedTitle}
\setbeamercolor{frametitle}{bg=VeryLightGray,fg=DarkText}
\setbeamercolor{normal text}{fg=DarkText}
\setbeamercolor{alerted text}{fg=RedTitle}
\setbeamersize{text margin left=5mm, text margin right=5mm}
\addtobeamertemplate{frametitle}{}{\vspace{-0.5em}}

\setbeamertemplate{navigation symbols}{}
\setbeamertemplate{footline}{%
  \hspace{5mm}{\usebeamerfont{page number in head/foot}\color{black!45}%
  ECON 282E \,$\cdot$\, \insertshortdeck}\hfill%
  \usebeamercolor[fg]{page number in head/foot}%
  \usebeamerfont{page number in head/foot}%
  \insertframenumber\,/\,\inserttotalframenumber\kern1em\vskip2pt%
}

% ---------------------------------------------------------------------------
% Chinese text under pdflatex (book titles etc.)
\newcommand{\zh}[1]{\begin{CJK*}{UTF8}{gbsn}#1\end{CJK*}}

% ---------------------------------------------------------------------------
% Deck title block
%   \decktitle{session}{deck id}{title}{date}
\newcommand{\insertshortdeck}{}
\newcommand{\decksession}{}
\newcommand{\deckid}{}
\newcommand{\decktitletext}{}
\newcommand{\deckdate}{}
\newcommand{\decktitle}[4]{%
  \renewcommand{\decksession}{#1}%
  \renewcommand{\deckid}{#2}%
  \renewcommand{\decktitletext}{#3}%
  \renewcommand{\deckdate}{#4}%
  \renewcommand{\insertshortdeck}{Session #1 \,$\cdot$\, Deck #1.#2}%
  \title{#3}%
  \author{Zhigang Feng}%
  \date{#4}%
}
\newcommand{\makedecktitle}{%
  \begin{frame}[plain,noframenumbering]
    \begin{tikzpicture}[remember picture,overlay]
      \fill[UCRBlue] (current page.north west) rectangle ([yshift=-0.55cm]current page.north east);
      \fill[UCRGold] ([yshift=-0.55cm]current page.north west) rectangle ([yshift=-0.62cm]current page.north east);
      \node[anchor=west, text=white, font=\small\bfseries] at ([xshift=5mm,yshift=-0.275cm]current page.north west)
        {ECON 282E \,$\cdot$\, Foundations of Macroeconomics};
      \node[anchor=east, text=white, font=\small] at ([xshift=-5mm,yshift=-0.275cm]current page.north east)
        {UC Riverside \,$\cdot$\, Fall 2026};
    \end{tikzpicture}
    \vspace*{1.1cm}
    {\usebeamerfont{title}\color{black!55}\large Session \decksession\ \,$\cdot$\, Deck \decksession.\deckid\par}
    \vspace{0.25cm}
    {\usebeamerfont{title}\color{RedTitle}\LARGE\bfseries \decktitletext\par}
    \vspace{0.7cm}
    {\large Zhigang Feng\par}
    \vspace{0.15cm}
    {\color{black!65}\deckdate\par}
    \vfill
    {\footnotesize\itshape\color{black!55}Build it on paper $\;\to\;$ solve it on the machine $\;\to\;$ push it to the frontier with AI\par}
  \end{frame}}

% ---------------------------------------------------------------------------
% Section divider (plain frame, not counted)
\newcommand{\sectiondivider}[2]{%
\begin{frame}[plain,noframenumbering]
\begin{center}
\vfill
{\Large\textbf{#1}\par}
\vspace{0.35cm}
{\normalsize #2\par}
\vfill
\end{center}
\end{frame}
}

% ---------------------------------------------------------------------------
% Boxes
\newtcolorbox{conceptbox}[1][]{colback=blue!3, colframe=RedTitle!55, boxrule=0.35mm,
  arc=1mm, left=2mm, right=2mm, top=1mm, bottom=1mm, #1}
\newtcolorbox{cautionbox}[1][]{colback=red!3, colframe=red!50!black, boxrule=0.35mm,
  arc=1mm, left=2mm, right=2mm, top=1mm, bottom=1mm, #1}
\newtcolorbox{examplebox}[1][]{colback=green!3, colframe=green!45!black, boxrule=0.35mm,
  arc=1mm, left=2mm, right=2mm, top=1mm, bottom=1mm, #1}
\newtcolorbox{takeawaybox}[1][]{colback=VeryLightGray, colframe=RedTitle!55, boxrule=0.35mm,
  arc=1mm, left=2mm, right=2mm, top=1mm, bottom=1mm, #1}
\newtcolorbox{researchbox}[1][]{enhanced, colback=violet!4, colframe=violet!60!black,
  boxrule=0.35mm, arc=1mm, left=2mm, right=2mm, top=1mm, bottom=1mm,
  fonttitle=\bfseries\small, title={Research in the wild}, #1}

\newcommand{\compacttables}{\renewcommand{\arraystretch}{1.15}}
\newcommand{\src}[1]{{\tiny\color{black!55}#1}}

% ---------------------------------------------------------------------------
\input{../preamble/course_map.tex}

%%   \decktitle{1}{A1}{Who I Am \& Why This Course}{September 24, 2026}
%%   \begin{document}
%%   \makedecktitle
%%   ...
%%
%% Adapted from Courses/AI-ECON-2026/source/shared_preamble.tex (Metropolis theme,
%% concept/caution/example/takeaway boxes, \sectiondivider). Changes for this course:
%%   * course title block (\decktitle / \makedecktitle) instead of \makebeamertitle
%%   * \zh{...} typesets Chinese (book titles) under pdflatex via CJKutf8
%%   * researchbox for the "Research in the wild" frames
%%   * section pages off; TikZ libraries and the course map loaded here
%% Build with pdflatex only (two passes); tools/build_slides.py does this for you.

\documentclass[10pt,english,aspectratio=169]{beamer}

\usepackage[T1]{fontenc}
\usepackage{textcomp}
\usepackage[utf8]{inputenc}
\usepackage{babel}
\usepackage{CJKutf8}

\usepackage{amsmath, amssymb, amsthm, graphicx}
\usepackage{booktabs, multirow, tabularx}
\usepackage{tikz}
\usetikzlibrary{positioning,arrows.meta,shapes,calc,fit,backgrounds}
\usepackage{pgfplots}
\pgfplotsset{compat=1.16}
\usepackage[most]{tcolorbox}
\tcbuselibrary{skins, breakable}
\usepackage{listings}
\usepackage{xcolor}

\lstset{
  basicstyle=\ttfamily\small,
  keywordstyle=\color{blue!80!black}\bfseries,
  commentstyle=\color{green!50!black}\itshape,
  stringstyle=\color{orange!80!black},
  backgroundcolor=\color{gray!8},
  frame=single,
  rulecolor=\color{gray!40},
  breaklines=true,
  showstringspaces=false,
  numbers=none,
}

\usetheme{metropolis}
\usefonttheme{professionalfonts}
\metroset{sectionpage=none}

\definecolor{LightGray}{RGB}{230,230,230}
\definecolor{RedTitle}{RGB}{0,0,200}
\definecolor{VeryLightGray}{RGB}{245,245,245}
\definecolor{DarkText}{RGB}{0,0,0}
\definecolor{UCRBlue}{RGB}{0,45,106}
\definecolor{UCRGold}{RGB}{241,171,0}

% Stage colours of the course arc (used by the course map and elsewhere)
\colorlet{StageTrunk}{blue!12}
\colorlet{StageBranch}{cyan!14}
\colorlet{StageMachine}{gray!18}
\colorlet{StageWall}{red!14}
\colorlet{StageTools}{orange!20}
\colorlet{StageData}{green!16}
\colorlet{StageAgents}{violet!14}

\hypersetup{bookmarks=false,breaklinks=true,pdfborder={0 0 0},colorlinks=true,
  linkcolor=DarkText,urlcolor=RedTitle}

\setbeamercolor{background canvas}{bg=white}
\setbeamercolor{title}{fg=RedTitle}
\setbeamercolor{frametitle}{bg=VeryLightGray,fg=DarkText}
\setbeamercolor{normal text}{fg=DarkText}
\setbeamercolor{alerted text}{fg=RedTitle}
\setbeamersize{text margin left=5mm, text margin right=5mm}
\addtobeamertemplate{frametitle}{}{\vspace{-0.5em}}

\setbeamertemplate{navigation symbols}{}
\setbeamertemplate{footline}{%
  \hspace{5mm}{\usebeamerfont{page number in head/foot}\color{black!45}%
  ECON 282E \,$\cdot$\, \insertshortdeck}\hfill%
  \usebeamercolor[fg]{page number in head/foot}%
  \usebeamerfont{page number in head/foot}%
  \insertframenumber\,/\,\inserttotalframenumber\kern1em\vskip2pt%
}

% ---------------------------------------------------------------------------
% Chinese text under pdflatex (book titles etc.)
\newcommand{\zh}[1]{\begin{CJK*}{UTF8}{gbsn}#1\end{CJK*}}

% ---------------------------------------------------------------------------
% Deck title block
%   \decktitle{session}{deck id}{title}{date}
\newcommand{\insertshortdeck}{}
\newcommand{\decksession}{}
\newcommand{\deckid}{}
\newcommand{\decktitletext}{}
\newcommand{\deckdate}{}
\newcommand{\decktitle}[4]{%
  \renewcommand{\decksession}{#1}%
  \renewcommand{\deckid}{#2}%
  \renewcommand{\decktitletext}{#3}%
  \renewcommand{\deckdate}{#4}%
  \renewcommand{\insertshortdeck}{Session #1 \,$\cdot$\, Deck #1.#2}%
  \title{#3}%
  \author{Zhigang Feng}%
  \date{#4}%
}
\newcommand{\makedecktitle}{%
  \begin{frame}[plain,noframenumbering]
    \begin{tikzpicture}[remember picture,overlay]
      \fill[UCRBlue] (current page.north west) rectangle ([yshift=-0.55cm]current page.north east);
      \fill[UCRGold] ([yshift=-0.55cm]current page.north west) rectangle ([yshift=-0.62cm]current page.north east);
      \node[anchor=west, text=white, font=\small\bfseries] at ([xshift=5mm,yshift=-0.275cm]current page.north west)
        {ECON 282E \,$\cdot$\, Foundations of Macroeconomics};
      \node[anchor=east, text=white, font=\small] at ([xshift=-5mm,yshift=-0.275cm]current page.north east)
        {UC Riverside \,$\cdot$\, Fall 2026};
    \end{tikzpicture}
    \vspace*{1.1cm}
    {\usebeamerfont{title}\color{black!55}\large Session \decksession\ \,$\cdot$\, Deck \decksession.\deckid\par}
    \vspace{0.25cm}
    {\usebeamerfont{title}\color{RedTitle}\LARGE\bfseries \decktitletext\par}
    \vspace{0.7cm}
    {\large Zhigang Feng\par}
    \vspace{0.15cm}
    {\color{black!65}\deckdate\par}
    \vfill
    {\footnotesize\itshape\color{black!55}Build it on paper $\;\to\;$ solve it on the machine $\;\to\;$ push it to the frontier with AI\par}
  \end{frame}}

% ---------------------------------------------------------------------------
% Section divider (plain frame, not counted)
\newcommand{\sectiondivider}[2]{%
\begin{frame}[plain,noframenumbering]
\begin{center}
\vfill
{\Large\textbf{#1}\par}
\vspace{0.35cm}
{\normalsize #2\par}
\vfill
\end{center}
\end{frame}
}

% ---------------------------------------------------------------------------
% Boxes
\newtcolorbox{conceptbox}[1][]{colback=blue!3, colframe=RedTitle!55, boxrule=0.35mm,
  arc=1mm, left=2mm, right=2mm, top=1mm, bottom=1mm, #1}
\newtcolorbox{cautionbox}[1][]{colback=red!3, colframe=red!50!black, boxrule=0.35mm,
  arc=1mm, left=2mm, right=2mm, top=1mm, bottom=1mm, #1}
\newtcolorbox{examplebox}[1][]{colback=green!3, colframe=green!45!black, boxrule=0.35mm,
  arc=1mm, left=2mm, right=2mm, top=1mm, bottom=1mm, #1}
\newtcolorbox{takeawaybox}[1][]{colback=VeryLightGray, colframe=RedTitle!55, boxrule=0.35mm,
  arc=1mm, left=2mm, right=2mm, top=1mm, bottom=1mm, #1}
\newtcolorbox{researchbox}[1][]{enhanced, colback=violet!4, colframe=violet!60!black,
  boxrule=0.35mm, arc=1mm, left=2mm, right=2mm, top=1mm, bottom=1mm,
  fonttitle=\bfseries\small, title={Research in the wild}, #1}

\newcommand{\compacttables}{\renewcommand{\arraystretch}{1.15}}
\newcommand{\src}[1]{{\tiny\color{black!55}#1}}

% ---------------------------------------------------------------------------
%% Course map: the recurring "you are here" slide for ECON 282E.
%% Loaded by course_preamble.tex. Pattern follows AI-ECON-2026 lec01_map.tex, but the map
%% shows the whole ten-session course rather than one lecture.
%%
%%   \CourseMap{s}{label}   s = 1..10 highlights session s and prints "you are here: label"
%%                          (e.g. \CourseMap{1}{1.A2}); earlier sessions get a check mark.
%%   \CourseMap{0}{}        full overview, nothing highlighted.
%%
%% Note: TikZ option lists are comma-split before TeX conditionals run, so every \ifnum
%% below sits outside [...] option lists.

\newcommand{\CourseMap}[2]{%
\begin{center}
\begin{tikzpicture}[
    sess/.style={rectangle, rounded corners=2pt, draw=black!45, minimum width=1.34cm,
        minimum height=1.12cm, text width=1.26cm, align=center, inner sep=1pt},
    stage/.style={font=\tiny\bfseries, text=black!70},
]
  \foreach \i/\d/\t/\c in {%
      1/Sep 24/Intro \\ Growth/StageTrunk,
      2/Oct 1/HA $\cdot$ OLG \\ Policy/StageBranch,
      3/Oct 8/Num.\ DP \\ Python/StageMachine,
      4/Oct 15/Numerics \\ I/StageMachine,
      5/Oct 22/Numerics \\ II/StageMachine,
      6/Oct 29/The wall \\ \textbf{Midterm}/StageWall,
      7/Nov 5/ML \\ Deep nets/StageTools,
      8/Nov 12/RL \\ HA $+$ DL/StageTools,
      9/Nov 19/Text \\ LLMs/StageData,
      10/Dec 3/Agents \\ \textbf{Projects}/StageAgents} {
    \node[sess, fill=\c] (n\i) at ({1.46*(\i-1)}, 0)
      {{\scriptsize\bfseries S\i}\\[-1pt]{\tiny\color{black!60}\d}\\[1pt]{\tiny \t}};
    \ifnum\i<#1
      \node[font=\scriptsize, text=green!45!black] at ([xshift=-2.5mm,yshift=-2.2mm]n\i.north east) {$\checkmark$};
    \fi
    \ifnum\i=#1
      \draw[RedTitle, very thick, rounded corners=2pt]
        ([xshift=-0.6mm,yshift=-0.6mm]n\i.south west) rectangle ([xshift=0.6mm,yshift=0.6mm]n\i.north east);
      % keep the label inside the picture at both ends of the row
      \ifnum\i<3
        \node[font=\scriptsize\bfseries, text=RedTitle, anchor=south west] at ([yshift=1.2mm]n\i.north west)
          {$\blacktriangledown$\,you are here: #2};
      \else\ifnum\i>8
        \node[font=\scriptsize\bfseries, text=RedTitle, anchor=south east] at ([yshift=1.2mm]n\i.north east)
          {you are here: #2\,$\blacktriangledown$};
      \else
        \node[font=\scriptsize\bfseries, text=RedTitle, anchor=south] at ([yshift=1.2mm]n\i.north)
          {you are here: #2\,$\blacktriangledown$};
      \fi\fi
    \fi
  }
  % stage band under the sessions
  \foreach \a/\b/\lab/\c in {1/1/trunk/StageTrunk, 2/2/branches/StageBranch,
      3/5/the machine $\cdot$ classical methods/StageMachine, 6/6/the wall/StageWall,
      7/8/new tools/StageTools, 9/9/new data/StageData, 10/10/colleagues/StageAgents} {
    \fill[\c] ([yshift=-1.5mm]n\a.south west) rectangle ([yshift=-4.6mm]n\b.south east);
    \path ([yshift=-1.5mm]n\a.south west) -- ([yshift=-4.6mm]n\b.south east)
      node[midway, stage] {\lab};
  }
  \node[font=\footnotesize\itshape, text=black!70] at ([yshift=-9.5mm]$(n1.south)!0.5!(n10.south)$)
    {Build it on paper $\;\to\;$ solve it on the machine $\;\to\;$ push it to the frontier with AI};
\end{tikzpicture}
\end{center}
}


%%   \decktitle{1}{A1}{Who I Am \& Why This Course}{September 24, 2026}
%%   \begin{document}
%%   \makedecktitle
%%   ...
%%
%% Adapted from Courses/AI-ECON-2026/source/shared_preamble.tex (Metropolis theme,
%% concept/caution/example/takeaway boxes, \sectiondivider). Changes for this course:
%%   * course title block (\decktitle / \makedecktitle) instead of \makebeamertitle
%%   * \zh{...} typesets Chinese (book titles) under pdflatex via CJKutf8
%%   * researchbox for the "Research in the wild" frames
%%   * section pages off; TikZ libraries and the course map loaded here
%% Build with pdflatex only (two passes); tools/build_slides.py does this for you.

\documentclass[10pt,english,aspectratio=169]{beamer}

\usepackage[T1]{fontenc}
\usepackage{textcomp}
\usepackage[utf8]{inputenc}
\usepackage{babel}
\usepackage{CJKutf8}

\usepackage{amsmath, amssymb, amsthm, graphicx}
\usepackage{booktabs, multirow, tabularx}
\usepackage{tikz}
\usetikzlibrary{positioning,arrows.meta,shapes,calc,fit,backgrounds}
\usepackage{pgfplots}
\pgfplotsset{compat=1.16}
\usepackage[most]{tcolorbox}
\tcbuselibrary{skins, breakable}
\usepackage{listings}
\usepackage{xcolor}

\lstset{
  basicstyle=\ttfamily\small,
  keywordstyle=\color{blue!80!black}\bfseries,
  commentstyle=\color{green!50!black}\itshape,
  stringstyle=\color{orange!80!black},
  backgroundcolor=\color{gray!8},
  frame=single,
  rulecolor=\color{gray!40},
  breaklines=true,
  showstringspaces=false,
  numbers=none,
}

\usetheme{metropolis}
\usefonttheme{professionalfonts}
\metroset{sectionpage=none}

\definecolor{LightGray}{RGB}{230,230,230}
\definecolor{RedTitle}{RGB}{0,0,200}
\definecolor{VeryLightGray}{RGB}{245,245,245}
\definecolor{DarkText}{RGB}{0,0,0}
\definecolor{UCRBlue}{RGB}{0,45,106}
\definecolor{UCRGold}{RGB}{241,171,0}

% Stage colours of the course arc (used by the course map and elsewhere)
\colorlet{StageTrunk}{blue!12}
\colorlet{StageBranch}{cyan!14}
\colorlet{StageMachine}{gray!18}
\colorlet{StageWall}{red!14}
\colorlet{StageTools}{orange!20}
\colorlet{StageData}{green!16}
\colorlet{StageAgents}{violet!14}

\hypersetup{bookmarks=false,breaklinks=true,pdfborder={0 0 0},colorlinks=true,
  linkcolor=DarkText,urlcolor=RedTitle}

\setbeamercolor{background canvas}{bg=white}
\setbeamercolor{title}{fg=RedTitle}
\setbeamercolor{frametitle}{bg=VeryLightGray,fg=DarkText}
\setbeamercolor{normal text}{fg=DarkText}
\setbeamercolor{alerted text}{fg=RedTitle}
\setbeamersize{text margin left=5mm, text margin right=5mm}
\addtobeamertemplate{frametitle}{}{\vspace{-0.5em}}

\setbeamertemplate{navigation symbols}{}
\setbeamertemplate{footline}{%
  \hspace{5mm}{\usebeamerfont{page number in head/foot}\color{black!45}%
  ECON 282E \,$\cdot$\, \insertshortdeck}\hfill%
  \usebeamercolor[fg]{page number in head/foot}%
  \usebeamerfont{page number in head/foot}%
  \insertframenumber\,/\,\inserttotalframenumber\kern1em\vskip2pt%
}

% ---------------------------------------------------------------------------
% Chinese text under pdflatex (book titles etc.)
\newcommand{\zh}[1]{\begin{CJK*}{UTF8}{gbsn}#1\end{CJK*}}

% ---------------------------------------------------------------------------
% Deck title block
%   \decktitle{session}{deck id}{title}{date}
\newcommand{\insertshortdeck}{}
\newcommand{\decksession}{}
\newcommand{\deckid}{}
\newcommand{\decktitletext}{}
\newcommand{\deckdate}{}
\newcommand{\decktitle}[4]{%
  \renewcommand{\decksession}{#1}%
  \renewcommand{\deckid}{#2}%
  \renewcommand{\decktitletext}{#3}%
  \renewcommand{\deckdate}{#4}%
  \renewcommand{\insertshortdeck}{Session #1 \,$\cdot$\, Deck #1.#2}%
  \title{#3}%
  \author{Zhigang Feng}%
  \date{#4}%
}
\newcommand{\makedecktitle}{%
  \begin{frame}[plain,noframenumbering]
    \begin{tikzpicture}[remember picture,overlay]
      \fill[UCRBlue] (current page.north west) rectangle ([yshift=-0.55cm]current page.north east);
      \fill[UCRGold] ([yshift=-0.55cm]current page.north west) rectangle ([yshift=-0.62cm]current page.north east);
      \node[anchor=west, text=white, font=\small\bfseries] at ([xshift=5mm,yshift=-0.275cm]current page.north west)
        {ECON 282E \,$\cdot$\, Foundations of Macroeconomics};
      \node[anchor=east, text=white, font=\small] at ([xshift=-5mm,yshift=-0.275cm]current page.north east)
        {UC Riverside \,$\cdot$\, Fall 2026};
    \end{tikzpicture}
    \vspace*{1.1cm}
    {\usebeamerfont{title}\color{black!55}\large Session \decksession\ \,$\cdot$\, Deck \decksession.\deckid\par}
    \vspace{0.25cm}
    {\usebeamerfont{title}\color{RedTitle}\LARGE\bfseries \decktitletext\par}
    \vspace{0.7cm}
    {\large Zhigang Feng\par}
    \vspace{0.15cm}
    {\color{black!65}\deckdate\par}
    \vfill
    {\footnotesize\itshape\color{black!55}Build it on paper $\;\to\;$ solve it on the machine $\;\to\;$ push it to the frontier with AI\par}
  \end{frame}}

% ---------------------------------------------------------------------------
% Section divider (plain frame, not counted)
\newcommand{\sectiondivider}[2]{%
\begin{frame}[plain,noframenumbering]
\begin{center}
\vfill
{\Large\textbf{#1}\par}
\vspace{0.35cm}
{\normalsize #2\par}
\vfill
\end{center}
\end{frame}
}

% ---------------------------------------------------------------------------
% Boxes
\newtcolorbox{conceptbox}[1][]{colback=blue!3, colframe=RedTitle!55, boxrule=0.35mm,
  arc=1mm, left=2mm, right=2mm, top=1mm, bottom=1mm, #1}
\newtcolorbox{cautionbox}[1][]{colback=red!3, colframe=red!50!black, boxrule=0.35mm,
  arc=1mm, left=2mm, right=2mm, top=1mm, bottom=1mm, #1}
\newtcolorbox{examplebox}[1][]{colback=green!3, colframe=green!45!black, boxrule=0.35mm,
  arc=1mm, left=2mm, right=2mm, top=1mm, bottom=1mm, #1}
\newtcolorbox{takeawaybox}[1][]{colback=VeryLightGray, colframe=RedTitle!55, boxrule=0.35mm,
  arc=1mm, left=2mm, right=2mm, top=1mm, bottom=1mm, #1}
\newtcolorbox{researchbox}[1][]{enhanced, colback=violet!4, colframe=violet!60!black,
  boxrule=0.35mm, arc=1mm, left=2mm, right=2mm, top=1mm, bottom=1mm,
  fonttitle=\bfseries\small, title={Research in the wild}, #1}

\newcommand{\compacttables}{\renewcommand{\arraystretch}{1.15}}
\newcommand{\src}[1]{{\tiny\color{black!55}#1}}

% ---------------------------------------------------------------------------
%% Course map: the recurring "you are here" slide for ECON 282E.
%% Loaded by course_preamble.tex. Pattern follows AI-ECON-2026 lec01_map.tex, but the map
%% shows the whole ten-session course rather than one lecture.
%%
%%   \CourseMap{s}{label}   s = 1..10 highlights session s and prints "you are here: label"
%%                          (e.g. \CourseMap{1}{1.A2}); earlier sessions get a check mark.
%%   \CourseMap{0}{}        full overview, nothing highlighted.
%%
%% Note: TikZ option lists are comma-split before TeX conditionals run, so every \ifnum
%% below sits outside [...] option lists.

\newcommand{\CourseMap}[2]{%
\begin{center}
\begin{tikzpicture}[
    sess/.style={rectangle, rounded corners=2pt, draw=black!45, minimum width=1.34cm,
        minimum height=1.12cm, text width=1.26cm, align=center, inner sep=1pt},
    stage/.style={font=\tiny\bfseries, text=black!70},
]
  \foreach \i/\d/\t/\c in {%
      1/Sep 24/Intro \\ Growth/StageTrunk,
      2/Oct 1/HA $\cdot$ OLG \\ Policy/StageBranch,
      3/Oct 8/Num.\ DP \\ Python/StageMachine,
      4/Oct 15/Numerics \\ I/StageMachine,
      5/Oct 22/Numerics \\ II/StageMachine,
      6/Oct 29/The wall \\ \textbf{Midterm}/StageWall,
      7/Nov 5/ML \\ Deep nets/StageTools,
      8/Nov 12/RL \\ HA $+$ DL/StageTools,
      9/Nov 19/Text \\ LLMs/StageData,
      10/Dec 3/Agents \\ \textbf{Projects}/StageAgents} {
    \node[sess, fill=\c] (n\i) at ({1.46*(\i-1)}, 0)
      {{\scriptsize\bfseries S\i}\\[-1pt]{\tiny\color{black!60}\d}\\[1pt]{\tiny \t}};
    \ifnum\i<#1
      \node[font=\scriptsize, text=green!45!black] at ([xshift=-2.5mm,yshift=-2.2mm]n\i.north east) {$\checkmark$};
    \fi
    \ifnum\i=#1
      \draw[RedTitle, very thick, rounded corners=2pt]
        ([xshift=-0.6mm,yshift=-0.6mm]n\i.south west) rectangle ([xshift=0.6mm,yshift=0.6mm]n\i.north east);
      % keep the label inside the picture at both ends of the row
      \ifnum\i<3
        \node[font=\scriptsize\bfseries, text=RedTitle, anchor=south west] at ([yshift=1.2mm]n\i.north west)
          {$\blacktriangledown$\,you are here: #2};
      \else\ifnum\i>8
        \node[font=\scriptsize\bfseries, text=RedTitle, anchor=south east] at ([yshift=1.2mm]n\i.north east)
          {you are here: #2\,$\blacktriangledown$};
      \else
        \node[font=\scriptsize\bfseries, text=RedTitle, anchor=south] at ([yshift=1.2mm]n\i.north)
          {you are here: #2\,$\blacktriangledown$};
      \fi\fi
    \fi
  }
  % stage band under the sessions
  \foreach \a/\b/\lab/\c in {1/1/trunk/StageTrunk, 2/2/branches/StageBranch,
      3/5/the machine $\cdot$ classical methods/StageMachine, 6/6/the wall/StageWall,
      7/8/new tools/StageTools, 9/9/new data/StageData, 10/10/colleagues/StageAgents} {
    \fill[\c] ([yshift=-1.5mm]n\a.south west) rectangle ([yshift=-4.6mm]n\b.south east);
    \path ([yshift=-1.5mm]n\a.south west) -- ([yshift=-4.6mm]n\b.south east)
      node[midway, stage] {\lab};
  }
  \node[font=\footnotesize\itshape, text=black!70] at ([yshift=-9.5mm]$(n1.south)!0.5!(n10.south)$)
    {Build it on paper $\;\to\;$ solve it on the machine $\;\to\;$ push it to the frontier with AI};
\end{tikzpicture}
\end{center}
}


\graphicspath{{./figures/}}

\lstset{language=Python, basicstyle=\ttfamily\scriptsize, frame=none,
        backgroundcolor=\color{gray!8}, aboveskip=2pt, belowskip=2pt}

\decktitle{5}{B1}{Projection: Making a Residual Small}{Thursday, October 22, 2026}

\begin{document}

\makedecktitle

%=============================================================================
\begin{frame}{Where We Are}
\CourseMap{5}{5.B1}
\vspace{-1mm}
\begin{center}
\small \textbf{Block B:} \textbf{5.B1} the residual $\;\cdot\;$ \textbf{5.B2} choosing a basis $\;\cdot\;$ \textbf{5.B3} Chebyshev $\;\cdot\;$ \textbf{5.B4} finite elements $\;\cdot\;$ \textbf{5.B5} determining the coefficients $\;\cdot\;$ \textbf{5.B6} the worked example.
\end{center}
\end{frame}

%=============================================================================
\begin{frame}{The Other Road}
\begin{columns}[T]
\begin{column}{0.54\textwidth}
\small
Block A gave up \textbf{global validity} to buy speed and dimension. We measured the cost: nothing at the steady state, $1.8$ decades at half of it, and everything at a kink.
\medskip

Block B makes the opposite trade. Keep the solution valid on the \textbf{whole domain}, constraints included --- but stop storing one number per grid point.
\medskip

The unknown becomes a short vector of \textbf{coefficients}, and ``solved'' is redefined: not ``the value function stopped moving'', and not ``the derivatives match at a point'', but
\[
  \textbf{the residual is zero.}
\]
\end{column}
\begin{column}{0.43\textwidth}
\begin{center}\footnotesize
\begin{tabularx}{\linewidth}{@{}>{\raggedright\arraybackslash}p{1.55cm} >{\raggedright\arraybackslash}X@{}}
\toprule
\textbf{Method} & \textbf{The unknown is} \\
\midrule
VFI (S3--S4) & $N$ values on a grid \\
perturbation & derivatives at one point \\
\textbf{projection} & \textbf{coefficients on a basis} \\
\bottomrule
\end{tabularx}
\end{center}
\medskip
\begin{examplebox}\footnotesize
And there is a fourth row waiting in S7: a neural network is this same move with a different, learned basis.
\end{examplebox}
\end{column}
\end{columns}
\end{frame}

%=============================================================================
\begin{frame}{The Model, With Labour This Time}
\begin{columns}[T]
\begin{column}{0.58\textwidth}
\small
\vspace{-2mm}
\begin{align*}
\text{(Euler)}\quad & \frac{1}{c_{t}}=\beta\mathbb{E}_{t}\!\left[\frac{1}{c_{t+1}}\Big(1+\alpha k_{t+1}^{\alpha-1}(e^{z_{t+1}}l_{t+1})^{1-\alpha}-\delta\Big)\right]\\
\text{(Labour)}\quad & \psi\frac{c_{t}}{1-l_{t}}=(1-\alpha)k_{t}^{\alpha}(e^{z_{t}}l_{t})^{1-\alpha}\frac{1}{l_{t}}\\
\text{(Resource)}\quad & c_{t}+k_{t+1}=k_{t}^{\alpha}(e^{z_{t}}l_{t})^{1-\alpha}+(1-\delta)k_{t}\\
\text{(Shock)}\quad & z_{t}=\rho z_{t-1}+\varepsilon_{t}
\end{align*}
The objective is unchanged from 5.A1: find $d$ with $H(d)=0$, where now
\[
  k_{t+1}=d_1(k_t,z_t),\qquad l_t=d_2(k_t,z_t).
\]
\end{column}
\begin{column}{0.39\textwidth}
\begin{cautionbox}[title=\textbf{Read the indices}]\footnotesize
Everything inside the $t+1$ expectation carries $t+1$ subscripts --- including the shock and labour inside the marginal product. Getting that wrong is the commonest transcription error in this literature, and it survives compilation.
\end{cautionbox}
\medskip
\begin{examplebox}\footnotesize
Endogenous labour is here to make the point that projection handles \emph{systems}: two unknown functions, two residuals, one coefficient vector.
\end{examplebox}
\end{column}
\end{columns}
\vspace{1mm}
\src{Quant\_Macro Lec.~9, ``Model Setup''; corrected against Fern\'andez-Villaverde \& Guerr\'on-Quintana, \emph{Solution Methods} V (Appendix), ``Equilibrium conditions''.}
\end{frame}

%=============================================================================
\begin{frame}{The Core Idea, in One Line}
\begin{columns}[T]
\begin{column}{0.55\textwidth}
\small
Restrict the unknown function to a finite-dimensional family:
\[
  d^{n}(x;\theta)\;=\;\sum_{i=0}^{n}\theta_{i}\,\psi_{i}(x)
\]
\begin{itemize}\small
  \item $\{\psi_i\}$ is a chosen \textbf{basis} --- polynomials, splines, orthogonal polynomials, tent functions.
  \item $\{\theta_i\}$ are the \textbf{unknown coefficients}, and there are finitely many.
\end{itemize}
\medskip
Substituting into the equilibrium conditions defines the \textbf{residual function}
\[
  R(x;\theta)\;=\;H\big(d^{n}(x;\theta)\big),
\]
which is what we try to make small.
\end{column}
\begin{column}{0.42\textwidth}
\begin{conceptbox}[title=\textbf{Two separate choices}]\small
\textbf{(1)} which basis $\{\psi_i\}$ --- 5.B2, 5.B3, 5.B4.\\[1mm]
\textbf{(2)} what ``small'' means --- 5.B5.
\medskip
Different answers to these two questions are the different names in the literature. The framework is one framework.
\end{conceptbox}
\medskip
\begin{examplebox}\footnotesize
Replace $\sum\theta_i\psi_i$ by $\mathcal{N}(x;\theta)$, a neural network, and the framework is unchanged. That is S7.
\end{examplebox}
\end{column}
\end{columns}
\vspace{1mm}
\src{Quant\_Macro Lec.~9, ``Projection: The Core Concept''.}
\end{frame}

%=============================================================================
\begin{frame}{Why the Name ``Projection''}
\begin{columns}[T]
\begin{column}{0.55\textwidth}
\small
The true $d$ lives in a big function space --- $C^{1}$, or $L^{2}$. Choosing a basis \textbf{projects} that infinite-dimensional problem onto a finite-dimensional subspace.
\medskip

We then cannot expect $H(d^{n})=0$ \emph{everywhere}; there are only $n+1$ coefficients. Instead we impose finitely many conditions:
\medskip
\begin{itemize}\small
  \item \textbf{Collocation:} $R(x_j;\theta)=0$ at $n+1$ chosen points.
  \item \textbf{Galerkin:} $\langle R(\cdot;\theta),\psi_i\rangle=0$ for each $i$ --- the residual is orthogonal to the basis.
  \item \textbf{Least squares:} minimise $\|R(\cdot;\theta)\|$ over the domain.
\end{itemize}
\medskip
Each is a sense in which the residual has been \emph{projected away}.
\end{column}
\begin{column}{0.42\textwidth}
\begin{takeawaybox}\footnotesize
The result is \textbf{global}: valid over the whole domain, not a neighbourhood --- provided the basis can represent the function and the projection is sensibly chosen.
\end{takeawaybox}
\medskip
\begin{examplebox}\footnotesize
Compare 5.A3: perturbation imposes conditions on \emph{derivatives at one point}. Projection imposes conditions on \emph{values over a domain}. Same residual, opposite discipline.
\end{examplebox}
\end{column}
\end{columns}
\vspace{1mm}
\src{Quant\_Macro Lec.~9, ``Why `Projection'? Mathematical Intuition''.}
\end{frame}

%=============================================================================
\begin{frame}{The Residual Is a Function --- Look at It}
\begin{columns}[T]
\begin{column}{0.53\textwidth}
\begin{center}
\begin{tikzpicture}
\begin{axis}[
    width=7.8cm, height=5.4cm,
    xmin=0.07, xmax=0.375, ymin=-0.045, ymax=0.105,
    xlabel={\scriptsize capital $k$}, ylabel={\scriptsize residual $R(k;\theta)$},
    tick label style={font=\scriptsize}, label style={font=\scriptsize},
    scaled ticks=false, /pgf/number format/fixed,
    /pgf/number format/precision=2,
    axis x line*=bottom, axis y line*=left,
    legend style={font=\tiny, draw=none, fill=none,
                  at={(0.98,0.97)}, anchor=north east}]
  \addplot[very thick, RedTitle] coordinates {
      (0.07481,0.09895) (0.07781,0.08891) (0.08080,0.07949) (0.08379,0.07065) (0.08678,0.06234) 
      (0.08978,0.05455) (0.09277,0.04723) (0.09576,0.04035) (0.09875,0.03389) (0.10175,0.02783) 
      (0.10474,0.02214) (0.10773,0.01680) (0.11072,0.01179) (0.11372,0.00709) (0.11671,0.00269) 
      (0.11970,-0.00143) (0.12269,-0.00529) (0.12569,-0.00889) (0.12868,-0.01225) 
      (0.13167,-0.01539) (0.13466,-0.01831) (0.13766,-0.02102) (0.14065,-0.02354) 
      (0.14364,-0.02587) (0.14663,-0.02802) (0.14963,-0.03001) (0.15262,-0.03183) 
      (0.15561,-0.03350) (0.15860,-0.03502) (0.16160,-0.03640) (0.16459,-0.03764) 
      (0.16758,-0.03876) (0.17057,-0.03975) (0.17357,-0.04063) (0.17656,-0.04139) 
      (0.17955,-0.04205) (0.18254,-0.04260) (0.18554,-0.04305) (0.18853,-0.04341) 
      (0.19152,-0.04368) (0.19451,-0.04386) (0.19751,-0.04396) (0.20050,-0.04398) 
      (0.20349,-0.04392) (0.20648,-0.04379) (0.20948,-0.04359) (0.21247,-0.04332) 
      (0.21546,-0.04299) (0.21845,-0.04259) (0.22145,-0.04214) (0.22444,-0.04162) 
      (0.22743,-0.04106) (0.23042,-0.04044) (0.23342,-0.03976) (0.23641,-0.03904) 
      (0.23940,-0.03828) (0.24239,-0.03746) (0.24539,-0.03661) (0.24838,-0.03571) 
      (0.25137,-0.03477) (0.25436,-0.03380) (0.25736,-0.03279) (0.26035,-0.03174) 
      (0.26334,-0.03066) (0.26633,-0.02955) (0.26933,-0.02840) (0.27232,-0.02723) 
      (0.27531,-0.02602) (0.27830,-0.02479) (0.28130,-0.02353) (0.28429,-0.02225) 
      (0.28728,-0.02094) (0.29027,-0.01961) (0.29327,-0.01826) (0.29626,-0.01688) 
      (0.29925,-0.01549) (0.30224,-0.01407) (0.30524,-0.01263) (0.30823,-0.01118) 
      (0.31122,-0.00971) (0.31421,-0.00822) (0.31721,-0.00672) (0.32020,-0.00520) 
      (0.32319,-0.00367) (0.32618,-0.00212) (0.32918,-0.00056) (0.33217,0.00102) (0.33516,0.00260) 
      (0.33815,0.00420) (0.34115,0.00581) (0.34414,0.00743) (0.34713,0.00906) (0.35012,0.01070) 
      (0.35312,0.01235) (0.35611,0.01401) (0.35910,0.01567) (0.36209,0.01735) (0.36509,0.01903) 
      (0.36808,0.02072) (0.37107,0.02241) (0.37406,0.02411) };
  \addlegendentry{2 coefficients}
  \addplot[very thick, orange!85!black] coordinates {
      (0.07481,0.00644) (0.07781,0.00432) (0.08080,0.00253) (0.08379,0.00103) (0.08678,-0.00022) 
      (0.08978,-0.00125) (0.09277,-0.00209) (0.09576,-0.00275) (0.09875,-0.00327) 
      (0.10175,-0.00365) (0.10474,-0.00393) (0.10773,-0.00410) (0.11072,-0.00419) 
      (0.11372,-0.00421) (0.11671,-0.00416) (0.11970,-0.00406) (0.12269,-0.00392) 
      (0.12569,-0.00374) (0.12868,-0.00353) (0.13167,-0.00329) (0.13466,-0.00304) 
      (0.13766,-0.00276) (0.14065,-0.00248) (0.14364,-0.00220) (0.14663,-0.00190) 
      (0.14963,-0.00161) (0.15262,-0.00132) (0.15561,-0.00103) (0.15860,-0.00075) 
      (0.16160,-0.00048) (0.16459,-0.00022) (0.16758,0.00003) (0.17057,0.00027) (0.17357,0.00050) 
      (0.17656,0.00071) (0.17955,0.00091) (0.18254,0.00110) (0.18554,0.00127) (0.18853,0.00142) 
      (0.19152,0.00156) (0.19451,0.00168) (0.19751,0.00179) (0.20050,0.00188) (0.20349,0.00196) 
      (0.20648,0.00202) (0.20948,0.00206) (0.21247,0.00209) (0.21546,0.00211) (0.21845,0.00211) 
      (0.22145,0.00210) (0.22444,0.00208) (0.22743,0.00204) (0.23042,0.00199) (0.23342,0.00193) 
      (0.23641,0.00186) (0.23940,0.00178) (0.24239,0.00169) (0.24539,0.00159) (0.24838,0.00149) 
      (0.25137,0.00137) (0.25436,0.00125) (0.25736,0.00113) (0.26035,0.00100) (0.26334,0.00086) 
      (0.26633,0.00072) (0.26933,0.00059) (0.27232,0.00044) (0.27531,0.00030) (0.27830,0.00016) 
      (0.28130,0.00002) (0.28429,-0.00012) (0.28728,-0.00026) (0.29027,-0.00039) (0.29327,-0.00052) 
      (0.29626,-0.00064) (0.29925,-0.00075) (0.30224,-0.00086) (0.30524,-0.00096) 
      (0.30823,-0.00105) (0.31122,-0.00113) (0.31421,-0.00120) (0.31721,-0.00126) 
      (0.32020,-0.00130) (0.32319,-0.00134) (0.32618,-0.00135) (0.32918,-0.00135) 
      (0.33217,-0.00134) (0.33516,-0.00130) (0.33815,-0.00125) (0.34115,-0.00118) 
      (0.34414,-0.00108) (0.34713,-0.00097) (0.35012,-0.00083) (0.35312,-0.00067) 
      (0.35611,-0.00049) (0.35910,-0.00028) (0.36209,-0.00005) (0.36509,0.00021) (0.36808,0.00050) 
      (0.37107,0.00082) (0.37406,0.00117) };
  \addlegendentry{4 coefficients}
  \addplot[very thick, black!50] coordinates {
      (0.07481,0.00006) (0.07781,-0.00000) (0.08080,-0.00003) (0.08379,-0.00005) (0.08678,-0.00005) 
      (0.08978,-0.00004) (0.09277,-0.00003) (0.09576,-0.00002) (0.09875,-0.00001) (0.10175,0.00001) 
      (0.10474,0.00002) (0.10773,0.00002) (0.11072,0.00003) (0.11372,0.00003) (0.11671,0.00003) 
      (0.11970,0.00003) (0.12269,0.00003) (0.12569,0.00003) (0.12868,0.00002) (0.13167,0.00002) 
      (0.13466,0.00001) (0.13766,0.00001) (0.14065,0.00000) (0.14364,-0.00000) (0.14663,-0.00001) 
      (0.14963,-0.00001) (0.15262,-0.00002) (0.15561,-0.00002) (0.15860,-0.00002) 
      (0.16160,-0.00002) (0.16459,-0.00002) (0.16758,-0.00002) (0.17057,-0.00002) 
      (0.17357,-0.00002) (0.17656,-0.00002) (0.17955,-0.00002) (0.18254,-0.00001) 
      (0.18554,-0.00001) (0.18853,-0.00001) (0.19152,-0.00000) (0.19451,-0.00000) (0.19751,0.00000) 
      (0.20050,0.00001) (0.20349,0.00001) (0.20648,0.00001) (0.20948,0.00001) (0.21247,0.00001) 
      (0.21546,0.00001) (0.21845,0.00002) (0.22145,0.00002) (0.22444,0.00002) (0.22743,0.00002) 
      (0.23042,0.00001) (0.23342,0.00001) (0.23641,0.00001) (0.23940,0.00001) (0.24239,0.00001) 
      (0.24539,0.00001) (0.24838,0.00000) (0.25137,0.00000) (0.25436,-0.00000) (0.25736,-0.00000) 
      (0.26035,-0.00000) (0.26334,-0.00001) (0.26633,-0.00001) (0.26933,-0.00001) 
      (0.27232,-0.00001) (0.27531,-0.00001) (0.27830,-0.00001) (0.28130,-0.00001) 
      (0.28429,-0.00001) (0.28728,-0.00001) (0.29027,-0.00001) (0.29327,-0.00001) 
      (0.29626,-0.00001) (0.29925,-0.00001) (0.30224,-0.00000) (0.30524,-0.00000) (0.30823,0.00000) 
      (0.31122,0.00000) (0.31421,0.00000) (0.31721,0.00001) (0.32020,0.00001) (0.32319,0.00001) 
      (0.32618,0.00001) (0.32918,0.00001) (0.33217,0.00001) (0.33516,0.00001) (0.33815,0.00001) 
      (0.34115,0.00001) (0.34414,0.00000) (0.34713,0.00000) (0.35012,-0.00000) (0.35312,-0.00000) 
      (0.35611,-0.00001) (0.35910,-0.00001) (0.36209,-0.00001) (0.36509,-0.00001) 
      (0.36808,-0.00001) (0.37107,-0.00000) (0.37406,0.00001) };
  \addlegendentry{8 coefficients}
  \draw[black!35] (axis cs:0.07,0) -- (axis cs:0.375,0);
  \addplot[only marks, mark=*, mark size=1.4pt, RedTitle, forget plot] coordinates {
      (0.33024,0.00000) (0.11864,0.00000) };
\end{axis}
\end{tikzpicture}
\end{center}
\end{column}
\begin{column}{0.44\textwidth}
\small
Brock--Mirman, deterministic, Chebyshev in $k$, coefficients chosen by collocation.
\medskip

\textbf{The residual is exactly zero at the collocation nodes} (dots, $2\times10^{-16}$) and free everywhere else. More coefficients push the whole curve towards the axis.
\medskip
\begin{center}\footnotesize
\begin{tabularx}{\linewidth}{@{}r r r@{}}
\toprule
coeff. & $\max|R|$ & policy error \\
\midrule
2 & $0.0990$ & $5.96\%$ \\
4 & $0.0064$ & $0.36\%$ \\
8 & $\mathbf{5.7\times10^{-5}}$ & $\mathbf{0.0038\%}$ \\
\bottomrule
\end{tabularx}
\end{center}
\end{column}
\end{columns}
\vspace{1mm}
\src{Ledger \texttt{projection\_residual}. Policy error is $\max|g/g^{*}-1|$ against the closed form $\alpha\beta k^{\alpha}$ on $[0.4k^{*},2k^{*}]$.}
\end{frame}

%=============================================================================
\begin{frame}{And Notice What the Residual Is Doing}
\begin{columns}[T]
\begin{column}{0.53\textwidth}
\small
In every row of that table the residual is \textbf{larger} than the policy error it is standing in for: $0.099$ against $0.060$, $0.0064$ against $0.0036$, $5.7\times10^{-5}$ against $3.8\times10^{-5}$.
\medskip

That is not a coincidence. It is the property 3.A1 needed from Santos (2000): the residual is \textbf{computable} and the policy error is \textbf{not}, and the theorem says the second is bounded by a constant times the first.
\medskip

Here, with the closed form in hand, we can check the bound directly instead of trusting it --- and the ratio sits between $1.5$ and $1.8$ at every basis size.
\end{column}
\begin{column}{0.44\textwidth}
\begin{takeawaybox}\footnotesize
So ``make the residual small'' is not a convenient proxy chosen for tractability. It is the only error measure you can compute, and there is a theorem licensing it.
\end{takeawaybox}
\medskip
\begin{cautionbox}\footnotesize
The constant depends on the curvature of $u$ and grows as the problem flattens. A residual is an \emph{order of magnitude}, not an error bar (3.A1).
\end{cautionbox}
\end{column}
\end{columns}
\vspace{1mm}
\src{Ledger \texttt{projection\_residual}; Santos (2000), and Santos \& Vigo-Aguiar (1998) for the $h^{2}/h$ pair.}
\end{frame}

%=============================================================================
\begin{frame}{The Algorithm, Formally}
\begin{columns}[T]
\begin{column}{0.57\textwidth}
\small
\begin{enumerate}\small\setlength{\itemsep}{1.2mm}
  \item Define $n+1$ linearly independent basis functions $\psi_i:\Psi\to\mathbb{R}^{m}$.
  \item Define the coefficient vector $\theta=[\theta_0,\dots,\theta_n]$.
  \item Form the approximant $\;d^{n}(\cdot;\theta)=\sum_{i=0}^{n}\theta_i\psi_i(\cdot)$.
  \item Substitute into $H$ to get the residual
  \[
    R(\cdot;\theta)=H\big(d^{n}(\cdot;\theta)\big).
  \]
  \item Choose $\hat{\theta}$ to make it small under a metric $\rho$:
  \[
    \hat{\theta}=\arg\min_{\theta}\;\rho\big[R(\cdot;\theta)\big].
  \]
\end{enumerate}
\end{column}
\begin{column}{0.40\textwidth}
\begin{conceptbox}[title=\textbf{Where the work is}]\small
Step 5 is a \textbf{nonlinear system} in $\theta$, solved by Newton or a quasi-Newton method --- so all of 4.A1 and 4.A3 apply here directly.
\end{conceptbox}
\end{column}
\end{columns}
\vspace{1mm}
\src{Quant\_Macro Lec.~9, ``Algorithm''.}
\end{frame}

%=============================================================================
\begin{frame}{Projection Against Perturbation}
\small
\begin{center}
\begin{tabularx}{0.97\linewidth}{@{}>{\raggedright\arraybackslash}p{2.2cm} >{\raggedright\arraybackslash}X >{\raggedright\arraybackslash}X@{}}
\toprule
 & \textbf{Perturbation (Block A)} & \textbf{Projection (Block B)} \\
\midrule
Approximates & locally, around one point & globally, on a domain \\
Unknowns are & derivatives at that point & coefficients on a basis \\
Conditions & derivatives of $F$ vanish & $R$ small at points, or in an inner product, or in norm \\
Speed & very fast; analytic derivatives & a nonlinear solve; slower \\
Dimension & tractable in large state spaces --- only local information is used & the curse of dimensionality returns \\
Nonlinearity & degrades away from the point & handled directly; no expansion \\
\bottomrule
\end{tabularx}
\end{center}
\medskip
\begin{center}
\begin{minipage}{0.86\linewidth}
\begin{conceptbox}[title=\textbf{Not rivals}]\footnotesize
They mix. Perturbation supplies the \textbf{initial guess} for the projection solve --- which is exactly what 5.B6 does, and why its nonlinear system converges at all.
\end{conceptbox}
\end{minipage}
\end{center}
\vspace{1mm}
\src{Quant\_Macro Lec.~9, ``Projection vs.\ Perturbation'' and ``Projection Method: Pros and Cons''.}
\end{frame}

%=============================================================================
\begin{frame}{Takeaways}
\begin{takeawaybox}
\begin{itemize}\small
  \item Projection replaces the unknown \emph{function} with a finite vector of \textbf{coefficients on a basis}, and redefines ``solved'' as \textbf{the residual is small}.
  \item Two independent choices: \textbf{which basis} (5.B2--5.B4), and \textbf{what ``small'' means} (5.B5). Every named method in the literature is one pair of answers.
  \item The residual is a \textbf{function}. Collocation zeroes it at $n+1$ points and leaves it free in between --- measured: $2\times10^{-16}$ at the nodes, $0.0990$ between them with two coefficients.
  \item More coefficients push the whole curve down: $0.0990 \to 0.0064 \to 5.7\times10^{-5}$ for $2$, $4$, $8$ coefficients, with policy error falling in step.
  \item The residual \textbf{exceeded} the policy error in every case, by a factor of $1.5$ to $1.8$ --- the property that makes Santos's theorem usable, checked here against a closed form rather than assumed.
  \item The solution is \textbf{global}: valid on the whole domain, kinks and constraints included, which is precisely where Block A had nothing to say.
\end{itemize}
\end{takeawaybox}
\end{frame}

%=============================================================================
\begin{frame}{Bridge: Which Basis?}
\begin{columns}[T]
\begin{column}{0.53\textwidth}
\small
Everything above was written with an unspecified $\{\psi_i\}$, and the example quietly used Chebyshev polynomials without saying why.
\medskip

The choice matters more than it looks. The obvious basis --- $1,x,x^{2},\dots$ --- fails, and it fails for a reason that has nothing to do with approximation theory and everything to do with conditioning.
\end{column}
\begin{column}{0.44\textwidth}
\begin{conceptbox}[title=\textbf{Next (5.B2)}]\small
Spectral against finite-element bases; why monomials are unusable; orthogonality; and the Gibbs phenomenon, which decides what happens when the function you are approximating has a kink.
\end{conceptbox}
\end{column}
\end{columns}
\end{frame}

\end{document}
